\eserciziario{Prototipo della fonte unica}

\section*{Problemi}

\livello{1}{Applicazione diretta delle formule}

\problema\label{pb:ud7-hvl-tessuto}
Un tessuto ha coefficiente di attenuazione \SI[per-mode=power]{0,35}{\centi\metre\tothe{-1}} per i raggi X di una radiografia. Calcola lo spessore emivalente.

\livello{2}{Problemi a due passaggi}

\problema\label{pb:ud7-piombo}
Per certi raggi $\gamma$ lo spessore emivalente del piombo è \SI{0,50}{\centi\metre}.
(a) Quale spessore riduce l'intensità a $1/16$?
(b) Quanti centimetri di piombo bisogna aggiungere ai precedenti perché passi meno dell'\SI{1}{\percent}, in numero intero di spessori emivalenti?

\livello{3}{Problemi di media difficoltà}

\problema\label{pb:ud7-tecnezio}
Per una scintigrafia un paziente deve ricevere alle 14:00 un radiofarmaco al tecnezio-99m (emivita \SI{6,0}{\hour}) con attività \SI{400}{\mega\becquerel}. La dose viene preparata alle 8:00 dello stesso giorno.
(a) Quale attività deve avere alla preparazione?
(b) Quanti nuclei di tecnezio-99m riceve il paziente?
(c) Trascurando l'eliminazione biologica, quale attività resta nel paziente 24 ore dopo la somministrazione?

\section*{Quesiti d'esame}

\quesito\label{qe:ud7-beta}
Il fosforo-32, usato in alcune terapie, decade $\beta^-$. Quale nucleo si forma?
\opzioni{$\nuclide{32}{14}{Si}$}{$\nuclide{32}{16}{S}$}{$\nuclide{31}{15}{P}$}{$\nuclide{28}{13}{Al}$}{$\nuclide{32}{15}{P}$}

\quesito\label{qe:ud7-assorbanza}
Una soluzione lascia passare l'\SI{1}{\percent} della luce che la colpisce. La sua assorbanza vale \completa.

\section*{Risultati rapidi}

\begin{center}
\begin{tabular}{@{}ll@{}}
\toprule
Esercizio & Risultato \\
\midrule
Problema~\ref{pb:ud7-hvl-tessuto} & \SI{2,0}{\centi\metre} \\
Problema~\ref{pb:ud7-piombo} & (a) \SI{2,0}{\centi\metre}; (b) \SI{1,5}{\centi\metre} \\
Problema~\ref{pb:ud7-tecnezio} & (a) \SI{800}{\mega\becquerel}; (b) \num{1,3e13}; (c) \SI{25}{\mega\becquerel} \\
Quesito~\ref{qe:ud7-beta} & B \\
Quesito~\ref{qe:ud7-assorbanza} & 2 \\
Quesito lampo~\ref{ql:ud7-hvl} & D \\
\bottomrule
\end{tabular}
\end{center}

\section*{Soluzioni svolte}

\subsection*{Problemi}

\solproblema{pb:ud7-hvl-tessuto}
Lo spessore emivalente è quello per cui $e^{-\mu x} = 1/2$:
\[
  x_{1/2} = \frac{\ln 2}{\mu} = \frac{0{,}69}{\SI[per-mode=power]{0,35}{\centi\metre\tothe{-1}}} \approx \SI{2,0}{\centi\metre}
\]

\solproblema{pb:ud7-piombo}
(a) $1/16 = 1/2^4$: servono 4 spessori emivalenti, cioè $4 \cdot \SI{0,50}{\centi\metre} = \SI{2,0}{\centi\metre}$.

(b) Con 6 spessori emivalenti passa $1/2^6 = 1/64$, più dell'\SI{1}{\percent}; con 7 passa $1/128$, meno dell'\SI{1}{\percent}. Servono 7 spessori in tutto, cioè $7 \cdot \SI{0,50}{\centi\metre} = \SI{3,5}{\centi\metre}$: ai \SI{2,0}{\centi\metre} di prima se ne aggiungono \SI{1,5}{\centi\metre}.


\solproblema{pb:ud7-tecnezio}
(a) Dalle 8:00 alle 14:00 passano \SI{6}{\hour}, cioè un'emivita: l'attività si dimezza. Alla preparazione serve il doppio, \SI{800}{\mega\becquerel}.

(b) Con l'emivita in secondi, $\SI{6,0}{\hour} = \SI{21600}{\second} \approx \SI{2,2e4}{\second}$:
\[
  \lambda = \frac{\ln 2}{T_{1/2}} \approx \frac{0{,}69}{\SI{2,2e4}{\second}} \approx \SI[per-mode=power]{3,2e-5}{\second\tothe{-1}}
  \qquad
  N = \frac{\mathcal{A}}{\lambda} = \frac{\SI[per-mode=power]{4,0e8}{\second\tothe{-1}}}{\SI[per-mode=power]{3,2e-5}{\second\tothe{-1}}} = 1{,}25 \cdot 10^{8-(-5)} \approx \num{1,3e13}
\]

(c) \SI{24}{\hour} sono $24/6 = 4$ emivite: $\SI{400}{\mega\becquerel} / 2^4 = \SI{25}{\mega\becquerel}$. In realtà l'attività sarà ancora più bassa, perché una parte del radiofarmaco viene eliminata con le urine.


\subsection*{Quesiti d'esame}

\solquesito{qe:ud7-beta}
\textbf{Risposta B.} Nel decadimento $\beta^-$ un neutrone diventa protone ed esce un elettrone: il numero di massa resta 32, il numero atomico sale da 15 a 16, che è lo zolfo.
\begin{itemize}
  \item A) $Z$ diminuisce di 1: è il decadimento $\beta^+$, non $\beta^-$.
  \item C) Toglie un neutrone al nucleo: nessun decadimento $\beta$ cambia il numero di massa.
  \item D) $A - 4$ e $Z - 2$: è il decadimento $\alpha$.
  \item E) Il nucleo resta lo stesso: succede solo nel decadimento $\gamma$.
\end{itemize}

\solquesito{qe:ud7-assorbanza}
\textbf{Risposta: 2.} Se passa l'\SI{1}{\percent}, $I_0/I = 100 = 10^2$ e $A = \log_{10} 10^2 = 2$.

\subsection*{Quesiti lampo}

\soluzionelampo{ql:ud7-hvl}
\textbf{Risposta D.} Ridurre a $1/4 = 1/2^2$ vuol dire attraversare 2 spessori emivalenti, quindi $x_{1/2} = \SI{1}{\centi\metre}$. \SI{5}{\centi\metre} sono 5 spessori emivalenti: $I/I_0 = 1/2^5 = 1/32$.
\begin{itemize}
  \item A) Tratta l'attenuazione come una proporzione: «$1/4$ in \SI{2}{\centi\metre}, quindi $1/10$ in \SI{5}{\centi\metre}». L'attenuazione è esponenziale, non lineare.
  \item B) È la frazione dopo \SI{4}{\centi\metre}: manca l'ultimo centimetro.
  \item C) Applica la legge dell'inverso del quadrato, che dipende dalla distanza dalla sorgente, non dal materiale.
  \item E) Scambia $1/4$ con $1/2$ e prende \SI{2}{\centi\metre} come spessore emivalente: $1/2^{2,5} \approx 1/6$.
\end{itemize}
